% ###################################################################################################
% Chapter 5 - Background
% ###################################################################################################
\chapter{Background}

The background section is designed to help players craft a rich and meaningful history for their characters. This section offers prompts to guide the player through the process of writing a cohesive backstory that connects the character's past with the decisions reflected on the character sheet (such as religion, scars, or personality traits). Below are questions that the player can use to generate a detailed narrative. Each answer should be written as part of a continuous text, creating a personal story for the character.

\subsection*{Birth and Early Life}

\begin{itemize}
    \item Where was your character born? \\
    \textit{Example: ``[Character Name] was born in [City Name], a small village nestled between mountains.''}
    
    \item What was their family situation like (parents, siblings, guardians)?
    \item Did they grow up in wealth or poverty?
\end{itemize}

\subsection*{Education and Early Influences}

\begin{itemize}
    \item What was their education like? Were they formally educated, self-taught, or completely untrained?
    \item Who were the people that had the most significant influence on them growing up?
    \item Did the character have a role model or mentor?
    \item What major events shaped the character’s worldview (war, loss, triumph)?
    \item Have they ever made a life-changing decision they regret or are proud of?
    \item Why did they adopt their religious or philosophical beliefs? What experiences led to this?
\end{itemize}

\subsection*{Appearance and Scars}

\begin{itemize}
    \item Are there any visible signs of past events on their body (scars, tattoos, piercings)?
    \item If so, what is the story behind each of these? \\
    \textit{Example: ``A long scar runs down his arm, a reminder of his first failed duel.''}
\end{itemize}

\subsection*{Relationships and Conflicts}

\begin{itemize}
    \item What relationships have shaped them (romantic, familial, friendships, rivalries)?
    \item Do they have any enemies or unresolved grudges?
    \item How do they handle conflicts and disagreements?
\end{itemize}

\subsection*{Occupation and Skills}

\begin{itemize}
    \item How did the character acquire their skills?
    \item What jobs or roles did they perform throughout their life?
\end{itemize}

\subsection*{Motivations and Goals}

\begin{itemize}
    \item What drives your character today? Do they have a primary goal or aspiration?
    \item Are they seeking redemption, revenge, recognition, or something else?
\end{itemize}

\subsection*{Secrets and Fears}

\begin{itemize}
    \item Does the character have any secrets? If so, what are they hiding and why?
    \item What are their greatest fears?
\end{itemize}

\subsection*{Encouraging Narrative Justification}

Players are encouraged to link their background story with other aspects of the character sheet, such as personality traits, religion, or scars. For example, if the character's religion reflects a personal experience, this should be included in the backstory. If they bear a particular scar, the background should explain how they got it and what impact it had on their life.

\subsection*{Player Tips}

\begin{itemize}
    \item \textbf{Focus on Motivation:} Every event in the character's life should build towards understanding what drives them today.
    \item \textbf{Leave Room for Growth:} Avoid writing a ``perfect'' character—leave flaws and conflicts unresolved to make space for in-game development.
    \item \textbf{Collaborate with the Game Master:} Share your background with the GM so they can weave elements into the campaign.
\end{itemize}

This approach will give players flexibility while offering enough structure to create meaningful and coherent character histories. The questions act as prompts to encourage creative storytelling while also ensuring that important character elements are tied together.

